% =========================================================
% Proof: Quotient Rule
% Source: volume-iii/analysis/differentiation/notes/notes-algebra.tex
% Mode: standard
% =========================================================

\subsection*{Quotient Rule}
\label{prf:quotient-rule}

\begin{remark*}[Return]
\hyperref[thm:quotient-rule]{$\leftarrow$ Back to Theorem (Quotient Rule) in Notes}
\end{remark*}

\begin{proof}
Let $I \subseteq \mathbb{R}$ be an open interval, let $c \in I$, and let
$f, g : I \to \mathbb{R}$ be differentiable at $c$ with $g(c) \neq 0$. Set
$h := f/g$. The claim is that $h$ is differentiable at $c$ and
\[
  h'(c) = \frac{f'(c)g(c) - f(c)g'(c)}{(g(c))^2}.
\]
Since $g$ is differentiable at $c$, it is continuous at $c$ by
\ByThm{Differentiability Implies Continuity}, so $\lim_{x \to c} g(x) = g(c)
\neq 0$. By \ByThm{Non-vanishing Limit}, $g(x) \neq 0$ for all $x$ in some
deleted neighbourhood of $c$, so $h = f/g$ is defined near $c$.

\ProofStep{1}{Form the difference quotient of $h$ and find a common denominator.}
For $x \in I$ with $x \neq c$ and $g(x) \neq 0$:
\[
  \frac{h(x) - h(c)}{x - c}
  = \frac{\dfrac{f(x)}{g(x)} - \dfrac{f(c)}{g(c)}}{x - c}
  = \frac{f(x)g(c) - f(c)g(x)}{g(x)g(c)(x - c)}.
\]

\ProofStep{2}{Insert the auxiliary term $f(c)g(c)$.}
Add and subtract $f(c)g(c)$ in the numerator:
\[
  \frac{f(x)g(c) - f(c)g(c) + f(c)g(c) - f(c)g(x)}{g(x)g(c)(x - c)}.
\]
Factor each pair and separate over the denominator:
\[
  = \frac{1}{g(x)g(c)}
    \left[
      g(c) \cdot \frac{f(x) - f(c)}{x - c}
      - f(c) \cdot \frac{g(x) - g(c)}{x - c}
    \right].
\]

\ProofStep{3}{Pass to the limit.}
By \ByThm{Sum Rule for Limits}, \ByThm{Product Rule for Limits}, and
\ByThm{Quotient Rule for Limits}:
\[
  \lim_{x \to c} \frac{h(x) - h(c)}{x - c}
  = \frac{1}{g(c) \cdot g(c)}
    \left[
      g(c) \cdot f'(c) - f(c) \cdot g'(c)
    \right],
\]
where $\lim_{x \to c} g(x) = g(c)$ by continuity of $g$, and the difference
quotient limits equal $f'(c)$ and $g'(c)$ by \ByDef{Differentiability at a
Point}. Therefore
\[
  h'(c) = \frac{g(c)f'(c) - f(c)g'(c)}{(g(c))^2}.
\]
\end{proof}

\begin{remark*}[Proof shape]
The proof mirrors the Product Rule: after combining fractions in the difference
quotient of $f/g$, the add-and-subtract trick $\pm f(c)g(c)$ decouples the
increments of $f$ and $g$ in the numerator, isolating the individual difference
quotients. Continuity of $g$ at $c$ supplies both $\lim g(x) = g(c)$ in the
denominator and — via the non-vanishing limit theorem — the guarantee that
$g(x) \neq 0$ near $c$ so the quotient is defined throughout.

\FlashProof{1. Combine the difference quotient of $f/g$ over a common
  denominator $g(x)g(c)(x-c)$.
  2. Insert $\pm f(c)g(c)$ in the numerator to isolate the difference
  quotients of $f$ and $g$.
  3. Apply limit laws; continuity of $g$ gives $\lim g(x) = g(c)$;
  differentiability gives $f'(c)$ and $g'(c)$.}
\end{remark*}

\begin{remark*}[Dependencies]
\begin{itemize}
  \item \hyperref[def:derivative-at-point]{Definition of Differentiability at a Point}:
    Identifies the difference quotient limits of $f$ and $g$ as $f'(c)$
    and $g'(c)$.
  \item \hyperref[thm:diff-implies-cont]{Differentiability Implies Continuity}:
    Supplies continuity of $g$ at $c$, giving $\lim_{x \to c} g(x) = g(c)$.
  \item \hyperref[thm:limit-nonvanishing]{Non-vanishing Limit}:
    Guarantees $g(x) \neq 0$ near $c$ so $f/g$ is defined in a neighbourhood
    of $c$.
  \item \hyperref[thm:limit-sum]{Sum Rule for Limits},
    \hyperref[thm:limit-product]{Product Rule for Limits},
    \hyperref[thm:limit-quotient]{Quotient Rule for Limits}:
    Distribute the limit across the algebraic expression in Step 3.
\end{itemize}

\FlashUses{def:derivative-at-point, thm:diff-implies-cont, thm:limit-nonvanishing, thm:limit-sum, thm:limit-product, thm:limit-quotient}
\end{remark*}